\documentclass[12pt]{article}

\usepackage{sbc-template}
\usepackage{graphicx,url}
\usepackage[utf8]{inputenc}
\usepackage[brazil]{babel}
\usepackage[latin1]{inputenc}


\sloppy

\title{Mecanismos de Backpressure com Apache Kafka}

\author{Willian Garbo G. Batista\inst{1}}


\address{Instituto de Informática -- Universidade Federal do ABC
  (UFABC)\\
  Caixa Postal 15.064 -- 91.501-970 -- São Paulo -- SP -- Brazil
  \email{wiilian.garbo@ufabc.edu.br}
}

\begin{document}

  \maketitle

  \begin{abstract}
    This meta-paper describes the style to be used in articles and short papers
    for SBC conferences. For papers in English, you should add just an abstract
    while for the papers in Portuguese, we also ask for an abstract in
    Portuguese (``resumo''). In both cases, abstracts should not have more than
    10 lines and must be in the first page of the paper.
  \end{abstract}

  \begin{resumo}
    Este meta-artigo descreve o estilo a ser usado na confecção de artigos e
    resumos de artigos para publicação nos anais das conferências organizadas
    pela SBC. É solicitada a escrita de resumo e abstract apenas para os artigos
    escritos em português. Artigos em inglês deverão apresentar apenas abstract.
    Nos dois casos, o autor deve tomar cuidado para que o resumo (e o abstract)
    não ultrapassem 10 linhas cada, sendo que ambos devem estar na primeira
    página do artigo.
  \end{resumo}


  \section{Introdução}

  Em sistemas distribuídos de alta volumetria de dados é essencial ter mecanismos que ajudem o sistemas a não sobrecarregarem em momentos de muito estresse. Filas, circuit breaker, rate limiting dentro outros são alguns dos mecanismos ou elementos que auxiliam os sistemas a lidarem de maneira que continuem em pleno funcionamento em picos de requisições, sendo eles exemplos de mecanismos de backpressure. No backpressure o desacoplamento entre os consumidores e produtores é um ponto muito importante, que é muitas vezes é trazido por filas ou brokens de mensageria sendo vital para consumidores processarem as mensagens de acordo com sua capacidade. Sem mecanismos de backpressure falhas em cascatas, instabilidades, e quedas podendo até acarretar em impossibilidades de uso. O Apache Kafka é uma ferramenta de processamento de streaming bastante utilizada em arquitetura orientada a eventos de alta volumetria, como em e-commerces e instituições financeiras, estando presente em mais de 80\% das empresas da Fortune 100  \cite{apachekafka:2026}, o que garantindo estabilidade e escalabilidade \cite{meissne:2018}. O Apache Kafka traz mecanismos de backpressure nativos  que ferramentas comuns não fornecem da mesma forma, como seu modelo de pub/sub baseado em pull, consumers group e partições tornando o robusto para lidar com até milhões de eventos por segundo.

  Este artigo apresenta uma avaliação experimental do Apache Kafka, permitindo o entendimento de alguns mecanismos nativos da ferramenta que auxiliam no backpressure em sistemas distribuidos. Dentro os mecanismos avaliados estão: Throttling/Quotas, Buffer, modelo de pubsub utilizando pull e consumer group.

  \section{Revisão da Literatura} \label{sec:firstpage}

  \subsection{Mecanismos de Backpressure}
  Em sistemas distribuídos, backpressure é um mecanismo de controle de fluxo utilizado para gerir o envio de dados de um produtor (upstream) para um consumidor (downstream) de forma que o dado produzido não sobrecarregue o consumidor, ou seja, que o dado seja transferido a uma taxa compatível com o processamento do consumidor \cite{fischer2022grokking, geeksforGeeks2026Jan}. A não utilização de mecanismos de backpressure pode causar desde aumento de latência, falhas de processamento e paralisação de sistemas, afetando diretamente os usuários finais dos sistemas. Além de afetar os custos de infraestrutura pela utilização excessiva de recursos nos consumidores (memória, cpu, network, filas etc), podendo ocorrer loops na inicialização da aplicação devido a sobrecarga que se se não monitorados causam danos financeiros altissimos \cite{newmanresilientds}. Os mecanismos são essenciais, pois sinalizam a produtores de dados, que os dados devem ser fornecidos a uma taxa menor ou devem ser paralizados por um período devido a sobrecarga dos consumidores. De acordo com Newman o backpressure pode ser implementado, do lado do cliente, como circuit break, e em ambos os lados, onde tanto o cliente quando o servidor entram em acordo, como rate limiting \cite{newmanresilientds}. A subseções abaixo explicam melhor sobre alguns mecanismos de backpressure.

  \subsubsection{Controle de Fluxo}
  É um mecanismo que faz a compatibilização de velocidade de produtor e consumidor \cite{kurose2013redes}, o que pode fazer com que os dados produzidos sejam recusados por um período até o consumidor consiga processar os dados. O controle de fluxo está presente no protocolo TCP, onde os dados são armazenados no buffer do consumidor, e de acordo com a  saturação do buffer o controle de fluxo passa a recursar mensagem do produtor. O produtor recebe uma sinalização do tamanho disponível do buffer do consumidor para evitar transmitir dados a uma velocidade que não poderá ser processada. Neste mecanismo de backpressure o produtor a partir de um dado enviado pelo consumidor consegue fazer a adaptação de sua taxa de transmissão para não sobrecarregar o consumidor e seu buffer.

  \subsubsection{Buffer}
  O buffer é uma estrutura de armazenamento de dados utilizada para guardar dados temporariamente para posterior processamento ou recuperação, por exemplo, uma fila. O buffer funciona como  backpressure, pois evita a sobrecarga direta do consumidor onde o mesmo pode recuperar os dados de acordo com sua velocidade de processamento. Contudo é necessário mecanimos de backpressure, como uma sinalização para o produtor informando para reduzir a velocidade de transmissão de dados ou cessá-la por um momento, para o buffer não chegar em seu limite máximo de armazenamento \cite{fischer2022grokking}.

  \subsubsection{Padrões de Arquitetura}
  \subsubsection{Rate limiting}
  É um padrão de arquitetura que restringe a quantidade de requisições por período de tempo (podendo ser em segundos, minutos ou horas) para cada cliente, deste modo, evitando consumo abusivo, o que poderia sobrecarregar o servidor e afetar o uso de recursos do servidor para outros clientes. Quando o limite de requisições é atingido o padrão passa a rejeitar novas requisições do cliente por uma mensagem (geralmente em HTTP status code 429 junto com uma mensagem no corpo). O padrão ajuda a evitar DoS (Denial of Service), assim mantendo a estabilidade do servidor. Neste padrão o backpressure está no fato de não permitir que um cliente consuma excessivamente o consumo de recursos de um servidor \cite{apiace:2023:europlop, agarwal2025Nov}.

  \subsubsection{Circuit Break}
  É um padrão de arquitetura onde a aplicação cliente passa a bloquear novas requisições ao servidor, quando um determinado limite de respostas com erros são retornadas pelo servidor acontece \cite{newmanresilientds}. Quando o limite é atingido, o circuito é aberto evitando novas solicitações por um determinado período, após esse período o circuito passa para o estado meio-aberto, onde passa a aceitar algumas requisições para validar se o problema continua, caso seja possivel fazer as requisições normalmente ao servidor o circuito é fechado, e o sistema volta ao estado normal. Caso contrário o mesmo volta para o estado de aberto. Esse padrão tem uma característica de falhar rápido o que permite o retorno de um feedback instântaneo sobre a requisição. Neste padrão o backpressue é implementado no lado cliente, onde o mesmo ativamente passa a evitar de enviar requisições quando existem erros retornados pelo servidor.

  \subsubsection{Load Shedding}
  É um padrão de arquitetura onde o sistema passa a rejeitar as requisições quando um determinado limite é atingido para se manter saudável \cite{newmanresilientds}, assim sendo possível continuar operando normalmente para algumas requisições. Quando o limite é atingido, novas solicitações não são processadas, e é dada uma mensagem ao cliente informando que não será possível tratar a requisição naquele momente (considerando uma API HTTP, poderia ser retornado o status code 429 Too Many Request ou 503 Service Unavailable). Esse é um outro mecanismo que atua como backpressure, pois quando um determinado limite é atingido o sistema passa a rejeitar algumas novas solicitações evitando sua sobrecarga.

  \subsection{Apache Kafka}
  Apache Kafka é uma ferramenta open-source distribuida de streaming de eventos focada em lidar com alta volumetria de dados com baixa latência e resiliência \cite{apacheKafka2026May, gorshkova2026kafka}. A plataforma atua como um broker de mensageria com alta disponibilidade e resiliente, sendo amplamente utilizada em sistemas críticos ou de alta volumetria que atuam com arquitetura orientada a eventos. Entre suas principais características estão seu funcionamento no modelo pubsub (publisher-subscriber), seu armazenamento imutável de eventos (append-only log) e sua capacidade de escalabilidade por sua funcionar em cluster, e por possuir partições de dados e consumer group.  Entre os conceitos essenciais da plataforma estão broker, produtor, tópicos, partições e consumidor \cite{confluenceKafka2026Jul}.

  \subsection{Mecanismos de Backpresseure com Apache Kafka}

  \subsubsection{Buffer}

  \subsubsection{Quotas/Throttling}

  \subsubsection{PubSub pull-based}

  \subsubsection{Consumer groups}

  \section{Metodologia}

  \section{Experimentos}

  \section{Conclusão}

  \bibliographystyle{sbc}
  \bibliography{sbc-template}

\end{document}